\scanpage{211}
\mspage{65}
\uline{Corollaire} 4.\deleted{\unsure{7}}8. Une flèche $u$ de $\mathrm{Diag}$ qui est une $\mathbb{D}$-équivalence universelle est $\mathbb{D}$-biasphérique, et a fortiori $\mathbb{D}$-connexe. Inversement, si $u$ est soit $\mathbb{D}$-lisse, soit $\mathbb{D}$-propre, soit une $\mathbb{D}$-fibration faible, et si $u$ est $\mathbb{D}$-biasphérique, \unsure{c'est} une $\mathbb{D}$-équiv. universelle.

La dém. est celle du \unsure{cas} particulier 4.3. On a \deleted{\unsure{montré d'ailleurs que}} \inserted{donc un diagramme} d'implications :

\diagram{\scriptsize
\begin{tikzcd}[ampersand replacement=\&,column sep=small,row sep=large]
\text{Cat-équivalence}\arrow[dr,Rightarrow] \&
 {\mathbb{D}\text{-équiv. universelle }(W_{\mathrm{univ}})}
 \arrow[d,Rightarrow,"{\substack{\text{implication stricte ; équivalence}\\
 \text{dans le cas où }u\text{ est soit }\mathbb{D}\text{-propre,}\\
 \text{soit }\mathbb{D}\text{-lisse, soit une }\mathbb{D}\text{-fibration faible}}}"] \& {}\\
{} \& {\mathbb{D}\text{-biasphérique}}\arrow[dl,Rightarrow]\arrow[d,Rightarrow,"{\substack{\text{équiv. ds}\\\text{tous les cas}\\\text{connus}}}"]\arrow[dr,Rightarrow] \& {}\\
{\mathbb{D}\text{-asphérique}}\arrow[dr,Rightarrow] \& {\mathbb{D}\text{-connexe}}\arrow[d,Rightarrow] \& {\mathbb{D}\text{-coasphérique}}\arrow[dl,Rightarrow]\\
{} \& {\mathbb{D}\text{-équivalence }\,(W)} \& {}
\end{tikzcd}}

% slanted
\margin{\uline{NB} \deleted{\illegible{1}} Les sept classes de morphismes de $\mathrm{Cat}$ sont stables par composition. Sauf la notion de $\mathbb{D}$-\unsure{connexité|connex.}, toutes ces notions ne dépendent que de l'ens $W=W_{\mathbb{D}}$ des $\mathbb{D}$-équivalences.}

Rappelons que pour $u\colon I'\mathbin{\unsure{-|\to}}I$ donné

a) $u$ est $\mathbb{D}$-asph. ssi $\forall$ \deleted{\unsure{$A\ (\in\mathbb{D}(e))$}}\inserted{$F$ \uline{constant} dans $\mathbb{D}(I)$}\unsure{,}

\[
 \tag*{\unsure{$(\alpha)$|$(a)$}}
 \deleted{\unsure{\xi_I|A_I}}F\xrightarrow{\overset{\mathrm{adj.}}{\sim}}
 u_*u^*\unsure{\{|(|\xi}\deleted{\unsure{\xi_I|A_I}}F
 \qquad\text{est iso.}
\]

b) $u$ est $\mathbb{D}$-coasph. sss $\forall$ \deleted{\unsure{$\xi$}}\unsure{$F$} \uline{\unsure{constant}}, dans $\mathbb{D}(I)$

\[
 \tag*{\unsure{$(\alpha)$|$(a)$|$(\beta)$}}
 \deleted{\unsure{\xi_I|A_I}}F\xleftarrow{\overset{\mathrm{adj}}{\sim}}
 u_!u^*(\deleted{\unsure{\xi_I|A_I}}F)\qquad\text{est iso.}
\]

c) $u$ est $\mathbb{D}$-\unsure{connexe|totalement} \deleted{\unsure{ssi}} sss pour \uuline{$\forall$} $F$ de $\mathbb{D}(I)$ \deleted{\unsure{qq}} (pas forcément constant, ni loc constant) \unsure{($\alpha$)|($a$)} est $\xrightarrow{\mathrm{iso}}$, $\underset{\unsure{\mathrm{iso}}}{\unsure{\leftarrow|\text{ou}|-}}$ ce qui revient au $=$, pour $\forall$ $F$ \unsure{($\alpha$)|($a$)}
% archive sheet number: 211
